% ============================================================
% METHODOLOGICAL REASSESSMENT MANUSCRIPT — QW-1660 + post-Phase16 audits
% ============================================================

\documentclass[11pt,a4paper]{article}

\usepackage{amsmath,amssymb,amsfonts}
\usepackage{graphicx}
\usepackage{geometry}
\usepackage{hyperref}
\usepackage{booktabs}
\usepackage{longtable}
\usepackage{float}
\usepackage{setspace}
\usepackage{authblk}
\usepackage{enumitem}

\geometry{margin=1in}
\onehalfspacing

\title{\textbf{Global Fractal Correlations in GW Noise:\\
Methodological Reassessment of QW-1660 (v5--v76)}}

\author[1]{Krzysztof Żuchowski}
\affil[1]{Independent Researcher, Fractal Information Theory Project}
\affil[ ]{\texttt{krzysztof.zuch@gmail.com}}
\date{2026-03-04}

\begin{document}
\maketitle

\begin{abstract}
This manuscript replaces the previous ``empirical validation'' wording with a strict methodological reassessment of the QW-1660 raw-strain program. We audit results from v5--v76 and later strict method checks (QW-1724, QW-1725, QW-2116), and separate robust findings from protocol-sensitive findings.

The central conclusion is conservative: \textbf{a stable global cross-detector fractal FIN correlation in raw LIGO/Virgo strain is not currently established at strict methodological level}. Several intermediate analyses reported nontrivial anti-persistent structure (e.g., cross-H values below 0.5), but later tests showed strong dependence on filtering, estimator choices, null construction, and protocol consistency. The current evidential status is \textbf{mixed/inconclusive}, with specific anomaly channels classified as non-robust in strict reanalysis.
\end{abstract}

\section{Scope and Data}
We analyze the already produced project artifacts, not new claims:
\begin{itemize}[leftmargin=1.5em]
    \item QW-1660 series (v5--v76) JSON outputs,
    \item strict GW method audits (QW-1724, QW-1725, QW-2116),
    \item Phase 76 ``true functional'' run (whitened H1-L1, MF-DFA q=0),
    \item external blind GW holdout branch (QW-2078, QW-2077) reported separately from raw-strain global-correlation claims.
\end{itemize}

\section{Protocol Inventory and Key Quantitative Points}
Table~\ref{tab:key_points} summarizes the critical numbers that drive the updated verdict.

\begin{longtable}{p{2.0cm}p{8.5cm}p{4.0cm}}
\caption{Key methodological checkpoints used in the reassessment}\label{tab:key_points}\\
\toprule
\textbf{Artifact} & \textbf{What it reports} & \textbf{Key value(s)} \\
\midrule
QW-1660 v5 & PSD scaling anomaly in selected setup & $Z=6.29$ \\
QW-1660 v13 & Cross-detector coherence check & $r=-0.145$, verdict: no coherence \\
QW-1660 v24 & Cross-Hurst vs shuffle/phase surrogates & H1-L1: $0.361$ vs $0.531/0.397$ \\
QW-1660 v26 & FIN multifractal summary & $H_0=0.2277$ (model-dependent stage) \\
QW-1660 v47 & Filter paradox test & unfiltered $0.0373$, filtered $0.1002$ \\
QW-1660 v57 & Null MC (simplified model) & obs $0.311$, ``5 sigma'' claim \\
QW-1660 v61 (strict rerun) & Full null with aligned protocol ($n=131072$, $200$ trials) & mean $0.3875$, obs $0.4086$, $p_{2s}=0.398$ \\
QW-1660 v63 (strict rerun) & Whitening test with null calibration ($n=131072$, $200$ trials) & whitened cross-H $0.7582$, $p_{2s}=0.00995$ \\
Phase 76 & Whitening test variant B (GWOSC fetch + MF-DFA q=0) & $H_{q=0}=0.5051$ \\
QW-1660 v65 (strict rerun) & Micro time-shift with permutation null ($n=131072$, $200$ perms) & best $|r|=0.0644$ at $-24.66$ ms, $p_{abs}=0.00498$; not near $10\pm2$ ms \\
QW-1724 & Method audit across v47--v65 & verdict: high-risk inconclusive \\
QW-1725 & Strict cross-H reanalysis & verdict: anomaly not robust \\
QW-2116 & Method-repair gate over v61/v63/v65 strict branch & verdict: method repair pass; scientific anomaly verdict unchanged (non-robust) \\
\bottomrule
\end{longtable}

\section{Methodological Findings}
\subsection{What appears reproducible}
\begin{itemize}[leftmargin=1.5em]
    \item Anti-persistent behavior ($H<0.5$) can be produced in multiple processing paths.
    \item Cross-detector statistics can depart from shuffle nulls in selected configurations.
\end{itemize}

\subsection{What fails strict robustness}
\begin{itemize}[leftmargin=1.5em]
    \item \textbf{Protocol drift and internal inconsistency:} QW-1724 identifies critical contradictions (projection inconsistency, scale-invariance failure, hard-coded observation risk in one null test branch).
    \item \textbf{Method-repair versus scientific claim:} QW-2116 confirms protocol repair (no hardcoded observation, aligned $n$, calibrated null/permutation tests), but the scientific branch verdict remains \texttt{GW\_CROSS\_HURST\_ANOMALY\_NOT\_ROBUST}.
    \item \textbf{Pair inconsistency:} strict reanalysis shows very different H1-L1 vs H1-V1 behavior under matched protocol, weakening ``global'' interpretation.
    \item \textbf{Lag inconsistency:} even when a local max-$|r|$ appears in strict rerun, it is off the expected light-travel window and therefore does not support a physically coherent broadband shared-phase signature.
    \item \textbf{Whitening dependence:} different whitening implementations produce substantially different outcomes (e.g., v63 vs Phase 76), indicating estimator/pipeline sensitivity.
\end{itemize}

\subsection{Strict synthesis}
Under strict criterion ``same conclusion under controlled protocol variants and nulls,'' the raw-strain global FIN-correlation claim is not closure-ready. Current best strict label is:
\begin{equation}
\texttt{GW\_RAW\_STRAIN\_GLOBAL\_FIN\_STATUS} = \texttt{INCONCLUSIVE / NOT\_ROBUST}
\end{equation}

\section{Relation to the External Blind GW Holdout Branch}
A separate branch (QW-2078 \textrightarrow{} QW-2077) reports threshold pass on a locked holdout feature protocol:
\[
\mathrm{AUC}=0.81496,\quad
\mathrm{ADV}^{q90}=0.31028,\quad
\mathrm{SEP}_{\mathrm{median}}=0.002056,\quad
\mathrm{GAP}_{\mathrm{median}}=0.001289.
\]
All preregistered thresholds pass in that branch.

This result is important, but it is \textbf{not identical} to proving a universal raw-strain global fractal background. It supports a specific locked prediction test, not full closure of the broader Phase-16 global-correlation claim.

\section{Updated Scientific Verdict}
\begin{enumerate}[leftmargin=1.5em]
    \item The prior wording ``empirical validation of global FIN background'' was too strong.
    \item The raw-strain QW-1660 line shows nontrivial effects but also major methodological fragility.
    \item QW-2116 confirms that methodology was repaired, but it does not upgrade the scientific anomaly verdict.
    \item Strict current status for global raw-strain FIN correlation is mixed/inconclusive.
    \item Claims must remain falsifiable and should be phrased as provisional until independent replication under locked protocol is completed.
\end{enumerate}

\section{Minimum Protocol for a Future Decisive Test}
\begin{itemize}[leftmargin=1.5em]
    \item Pre-register one primary estimator and one fallback estimator (no post-hoc switching).
    \item Lock GPS/event list and detector-pair selection before running.
    \item Lock whitening/filter pipeline and edge-cropping rules.
    \item Use multiple null families (shuffle, phase-randomized, PSD-matched synthetic, circular shifts) with fixed trial count.
    \item Publish complete lag profiles and pairwise consistency statistics.
    \item Require agreement across independent environments and independent teams.
\end{itemize}

\section*{Reproducibility Pointers}
\begin{itemize}[leftmargin=1.5em]
    \item QW-1660 artifacts: \texttt{QW\_1660\_v*.json}
    \item Strict audits: \texttt{report\_qw1724\_gw\_method\_audit.json}, \texttt{report\_qw1725\_gw\_strict\_cross\_hurst\_reanalysis.json}, \texttt{report\_qw2116\_gw1660\_method\_repair\_gate.json}
    \item Strict rerun scripts: \texttt{phase61\_full\_null\_model\_strict.py}, \texttt{phase63\_whitening\_test\_strict.py}, \texttt{phase65\_micro\_timeshift\_strict.py}
    \item Phase 76 artifacts: \texttt{FIN\_Search\_LIGO\_Phase76\_TrueFunctional.py}, \texttt{phase76\_true\_fin\_results.json}
\end{itemize}

\end{document}
